\documentclass[12pt,a4paper]{article}
\usepackage[margin=25mm, headheight=15pt, headsep=10pt]{geometry}
\usepackage{fontspec}
\usepackage[table]{xcolor} % Note: table option is required for rowcolors
\usepackage{titlesec}
\usepackage{tocloft}
\usepackage{fancyhdr}
\usepackage{booktabs}
\usepackage{tcolorbox}
\tcbuselibrary{skins, breakable}
\usepackage[colorlinks=true, linkcolor=emerald, urlcolor=emerald, bookmarks=true]{hyperref}

\definecolor{emerald}{RGB}{0,102,51}
\definecolor{darkred}{RGB}{139,0,0}
\definecolor{lightgray}{RGB}{245,245,245}

\usepackage[nil]{babel}
\babelprovide[import, main]{bengali}
\babelprovide[import]{arabic}
\babelfont{rm}[Renderer=HarfBuzz,Scale=1.0]{Noto Serif Bengali}
\babelfont{sf}[Renderer=HarfBuzz]{Noto Sans Bengali}
\babelfont{tt}[Renderer=HarfBuzz]{Noto Sans Bengali}
\babelfont[arabic]{rm}[Renderer=HarfBuzz,Scale=1.1]{Noto Naskh Arabic}

% Section styling
\titleformat{\section}{\Large\bfseries\color{emerald}}{\thesection.}{0.5em}{}
\titleformat{\subsection}{\large\bfseries\color{emerald!85!black}}{\thesubsection.}{0.5em}{}
\titleformat{\subsubsection}{\normalsize\bfseries\color{emerald!70!black}}{\thesubsubsection.}{0.5em}{}
\titlespacing*{\section}{0pt}{18pt}{8pt}

% Fancy Header/Footer setup
\pagestyle{fancy}
\fancyhf{} % clear all header and footer fields
\fancyhead[L]{\color{emerald}\small\textbf{\nouppercase{\leftmark}}}
\fancyfoot[L]{\scriptsize\color{black!50}{\lat{AI}}-সংকলিত, আলেম-যাচাইকৃত নয়}
\fancyfoot[C]{\color{emerald!70!black}\thepage}
\renewcommand{\headrulewidth}{0.4pt}
\renewcommand{\headrule}{\hbox to\headwidth{\color{emerald}\leaders\hrule height \headrulewidth\hfill}}
\renewcommand{\footrulewidth}{0pt}

% Custom boxes
% 1. Ayat/Hadith Box
\newtcolorbox{ayatbox}[1][]{
    enhanced, breakable, 
    colback=emerald!5, 
    colframe=emerald, 
    boxrule=0.5pt, 
    arc=3pt, 
    left=10pt, right=10pt, top=10pt, bottom=10pt,
    #1
}

% 2. Quote/Translation Box
\newtcolorbox{quotebox}[1][]{
    enhanced, breakable,
    frame hidden,
    colback=lightgray,
    borderline west={3pt}{0pt}{emerald},
    left=15pt, right=10pt, top=10pt, bottom=10pt,
    sharp corners,
    #1
}

% 3. AI-generated notice box
\newtcolorbox{warnbox}[1][]{
    enhanced, breakable,
    colback=red!4,
    colframe=darkred,
    boxrule=0.8pt,
    arc=3pt,
    left=10pt, right=10pt, top=10pt, bottom=10pt,
    #1
}

% Commands
\newcommand{\arab}[1]{\foreignlanguage{arabic}{#1}}
\newcommand{\kib}[1]{\textcolor{emerald}{\textbf{#1}}}

% Latin (English) fallback: Noto Serif Bengali has NO Latin glyphs
\newfontfamily\latfont[Renderer=HarfBuzz]{Noto Serif}
\newcommand{\lat}[1]{{\latfont #1}}

\setlength{\parskip}{5pt}
\setlength{\parindent}{0pt}

% Fix for missing bullet in Noto Serif Bengali
\renewcommand{\labelitemi}{$\bullet$}



\begin{document}
\begin{center}{\Large\bfseries\color{emerald} ঘটনা ৫ — ইবনে আব্বাস (রাঃ): নবীজির সঙ্গে তাহাজ্জুদ — খালার ঘরের রাত}\end{center}
\vspace{1em}
\section*{ঘটনা ৫ — ইবনে আব্বাস (রাঃ): তাহাজ্জুদে নবীজির সঙ্গে খালার ঘরের রাত}

\begin{warnbox}
 \textbf{গুরুত্বপূর্ণ নোটিশ (\lat{Important} \lat{Notice}):} এই কন্টেন্টটি \lat{AI} (কৃত্রিম বুদ্ধিমত্তা) দিয়ে প্রস্তুত ও সংকলিত। \lat{AI} কঠোর সূত্র-মান নীতি মেনে শুধু নির্ভরযোগ্য উৎস থেকে তথ্য সংগ্রহ করেছে — কেবল সহীহ/হাসান হাদিস ও সহীহ সনদের আসার রাখা হয়েছে, দুর্বল সনদ বাদ দেওয়া হয়েছে। তবে এটি কোনো আলেম বা মুহাদ্দিস কর্তৃক যাচাইকৃত নয়।
\end{warnbox}

\begin{quote}
\textbf{সম্পর্কিত ফাইল:} পার্ট ৫ঘ (তাহাজ্জুদ-নফল), পার্ট ৪ (পদ্ধতি)।
\end{quote}

\medskip\hrule\medskip

\subsection*{১. সূত্র কার্ড (\lat{Source} \lat{Card})}

\begin{table}[h]\centering\small
\begin{tabular}{ll}
বিষয় & বিবরণ \\
\hline
\textbf{ঘটনা} & কিশোর ইবনে আব্বাস (রাঃ) খালা মাইমুনার (রাঃ) ঘরে রাত কাটিয়ে নবীজির (সা.) তাহাজ্জুদ ও ওযু-পদ্ধতি প্রত্যক্ষ করেন \\
\textbf{রাবি} & আবদুল্লাহ ইবনে আব্বাস (রাঃ) \\
\textbf{গ্রন্থ ও নম্বর} & সহীহ বুখারি ১৩৮; সহীহ মুসলিম ৭৬৩ \\
\textbf{সনদের মান} & সহীহ \\
\hline
\end{tabular}
\end{table}

\medskip\hrule\medskip

\subsection*{২. পটভূমি (\lat{Background})}

\textbf{কে:} ইবনে আব্বাস (রাঃ) — নবীজির চাচাতো ভাই, তখন অল্প বয়স্ক কিশোর; তাঁর খালা মাইমুনা বিনতে হারিস (রাঃ) — নবীজির (সা.) স্ত্রী।
\textbf{কখন:} মদিনা যুগের মধ্যবর্তী সময় — নির্দিষ্ট তারিখ হয়নি।
\textbf{কোথায়:} মাইমুনা (রাঃ)-এর ঘর — যেখানে নবীজি (সা.) সেদিন রাতে অবস্থান করতেন।
\textbf{কেন:} ইবনে আব্বাস (রাঃ) নবীজির (সা.) রাতের ইবাদত-পদ্ধতি — ওযু, তাহাজ্জুদ, বিতর, শেষ-রাতের দুআ — নিজে শিখতে ও প্রত্যক্ষ করতে চেয়েছিলেন।

\medskip\hrule\medskip

\subsection*{৩. পূর্ণ ঘটনার বর্ণনা (\lat{The} \lat{Full} \lat{Event})}

ইবনে আব্বাস (রাঃ) বলেন:

\begin{quote}
\lat{“}নবীজি (সা.) যখন আমার খালা মাইমুনার (রাঃ) ঘরে রাত কাটাতেন, তখন আমি সেই ঘরে রাত কাটাতাম — নবীজির (সা.) নামাজ ও ইবাদত দেখার জন্য। রাতের গভীরে নবীজি (সা.) ঘুম থেকে জাগলেন; উঠে ওযু করার জন্য ঝুলন্ত মশকের (পানির থলি) কাছে গেলেন — দাঁড়িয়ে ওযু করলেন — এমনই নিখুঁত ওযু। তারপর তাহাজ্জুদে দাঁড়িয়ে গেলেন। আমি ওযু করার ইচ্ছায় এগিয়ে গেলাম — নবীজির (সা.) পাশে দাঁড়ালাম। নবীজি (সা.) বাম হাত আমার মাথায় রাখলেন, কান ধরে সামান্য টানলেন — আমি এগিয়ে গেলাম। নবীজি (সা.) তাকবির দিয়ে নামাজ শুরু করলেন — আমি তাঁর বাম পাশে দাঁড়ালাম। তিনি আমার ডান পাশে আমাকে দাঁড় করালেন — তারপর নামাজ পড়লেন। আমি তার পেছনে-পেছনে অনুসরণ করলাম।\lat{”}
\end{quote}

ইবনে আব্বাস (রাঃ) বর্ণনা করেন:

\begin{quote}
\lat{“}নবীজি (সা.) যখন রুকু করতেন — আমি তাঁর পাশে দাঁড়ানোর প্রস্তুতি নিই; তিনি আমাদের শিক্ষা দিয়ে বললেন: ... তিনি (নবীজি) সূরা আল-বাকারাহ, আল-ইমরান, আন-নিসা, আল-মায়িদাহ বা আল-আনআম — এই দীর্ঘ সূরাগুলো পড়লেন। ... অবশেষে তিনি (নবীজি) এক রাকাত বিতর পড়লেন। তারপর সালাম দিয়ে শেষ করেন — তারপর আমরা দুজনে (নবীজি ও ইবনে আব্বাস) লম্বা হয়ে শুয়ে পড়লাম।\lat{”} (সহীহ বুখারি ১৩৮; সহীহ মুসলিম ৭৬৩ — সংক্ষিপ্ত-সামগ্রিক)
\end{quote}

এই রাত থেকেই ইবনে আব্বাস (রাঃ) শিখলেন — নবীজির (সা.) ওযু, তিলাওয়াতের নিরবতা/সুস্পষ্টতা, তাহাজ্জুদের রাকাত-পদ্ধতি, বিতর, এবং শেষে গভীর প্রতিফলন। নবীজি (সা.) কিশোর ইবনে আব্বাসকে নিজের \textbf{ডান পাশে} দাঁড় করানোর মাধ্যমে ইমাম-মুক্তাদির সঠিক অবস্থান শেখালেন —

\begin{quote}
ইমামের সাথে নামাজ পড়ার সময় মুক্তাদি দাঁড়ায় \textbf{ইমামের ডানে}; এবং একাকী নামাজে না — (এই শিক্ষাই হাদিসে বর্ণিত)।
\end{quote}

\medskip\hrule\medskip

\subsection*{৪. ধাপে ধাপে বিশ্লেষণ (\lat{Step-by-Step} \lat{Analysis})}

\begin{enumerate}
\item \textbf{শিক্ষার্থীর দৃঢ় ইচ্ছা:} কিশোর হলেও নবীজির (সা.) ইবাদত দেখার জন্য রাত জেগে থাকা — ইলমের পিপাসা।
\item \textbf{ওযুর নিখুঁততা:} দাঁড়িয়ে-মশক থেকে ওযু — নবীজির (সা.) পূর্ণ ওযু-পদ্ধতির প্রত্যক্ষণ।
\item \textbf{নবীজির (সা.) স্নেহ-শিক্ষা:} হাত-কানে টেনে আদরে ডান পাশে দাঁড় করানো — শৃঙ্খলা-শিক্ষা কোমলভাবে।
\item \textbf{দীর্ঘ তিলাওয়াতের প্রামাণ্য:} রাতের নামাজে দীর্ঘ সূরা (বাকারাহ-ইমরান-নিসা-মায়িদাহ) — তাহাজ্জুদ-পদ্ধতির দ্বিতীয় প্রামাণ্য (হুজাইফার (রাঃ) বর্ণনার সাথে মিল)।
\item \textbf{বিতর-শিক্ষা:} তাহাজ্জুদ শেষে এক রাকাত বিতর — নবীজির (সা.) বিতরের ঘন ঘন আমল।
\item \textbf{ইমাম-মুক্তাদির অবস্থান:} মুক্তাদির ডানে দাঁড়ানোর সুন্নত-শিক্ষা — এখান থেকেই ফিকহ-বর্ণনা।
\end{enumerate}
\medskip\hrule\medskip

\subsection*{৫. শব্দের ব্যাখ্যা (\lat{Key} \lat{Words})}

\begin{itemize}
\item \textbf{তাহাজ্জুদ:} ঘুম ভেঙে উঠে রাতের নামাজ — নবীজির (সা.) নিয়মিত আমল।
\item \textbf{বিতর:} রাতের নামাজের শেষ — বিজোড় রাকাত (১/৩/৫...) — নবীজির (সা.) স্থায়ী আমল।
\item \textbf{মুক্তাদি}: ইমামের পেছনে নামাজ-অনুসরণকারী; ডানে দাঁড়ানো সুন্নত (একজন হলে)।
\end{itemize}
\medskip\hrule\medskip

\subsection*{৬. সংশ্লিষ্ট হাদিস ও আয়াত}

\begin{quote}
\textbf{\lat{“}রাতের কিছু অংশে তাহাজ্জুদ পড়ো — তোমার জন্য অতিরিক্ত ইবাদতস্বরূপ; আশা করা যায় তোমার রব তোমাকে প্রশংসিত স্থানে দাঁড় করাবেন।\lat{”}} — সূরা আল-ইসরা ১৭:৭৯ (বাংলা অর্থ)।
\end{quote}

\begin{quote}
\textbf{\lat{“}তারা রাতের অল্পই ঘুমাতো; আর প্রভাতে ক্ষমা প্রার্থনা করত।\lat{”}} — সূরা আল-যারিয়াত ৫১:১৭-১৮ (বাংলা অর্থ)।
\end{quote}

\begin{quote}
নবীজি (সা.): \textbf{\lat{“}বিতর তোমাদের উপর ওয়াজিব... (দৃঢ় তাগিদ)।\lat{”}} — (সহীহ বুখারি ৯৯০-এর আলোকে)।
\end{quote}

\medskip\hrule\medskip

\subsection*{৭. গভীর শিক্ষা (\lat{Deep} \lat{Lessons})}

\textbf{শিক্ষা ১ — ইলম-অন্বেষণে কোনো বয়স বাধা নয়:} কিশোর ইবনে আব্বাস (রাঃ) নবীজির (সা.) নিকটতম পর্যবেক্ষক হয়ে উঠলেন — পরে উম্মতের প্রথম মুফাসসির (कुरআন-ব্যাখ্যাকারী)।

\textbf{শিক্ষা ২ — শিক্ষকের কোমল পদ্ধতি:} আদর-সোহাগে শৃঙ্খলা-শিক্ষা — শিক্ষকতা (\lat{teaching}) এর সুন্দর উদাহরণ।

\textbf{শিক্ষা ৩ — তাহাজ্জুদ ও বিতরের আমল:} দীর্ঘ রাতের ইবাদতের মাঝে বিতর-শিক্ষা — নফল অনুশীলনের ধারাবাহিকতা।

\textbf{শিক্ষা ৪ — ইমাম-মুক্তাদির সম্পর্ক:} সঠিক অবস্থান ও অনুসরণের নিয়ম — জামাতের একতা এখান থেকেই।

\textbf{শিক্ষা ৫ — শিশুদের ইবাদতে উৎসাহ:} নবীজি (সা.) কিশোরকে নামাজে সাথে রেখেছেন — বালক-বালিকাদের নামাজ-অনুশীলনের গুরুত্ব।

\medskip\hrule\medskip

\subsection*{৮. জীবনে প্রয়োগের পদ্ধতি (\lat{Practical} \lat{Application})}

\begin{enumerate}
\item \textbf{শিক্ষার্থীর আগ্রহ:} নিজে অজানা বিষয় (নামাজের সূক্ষ্মতা) শেখার জন্য সময় করা — বয়স-বাধা ভাঙা।
\item \textbf{শিক্ষকতা (\lat{Parenting}/\lat{teaching}):} সন্তান/ছাত্রকে আদর দিয়ে নিয়ম শেখান — কঠোরতা নয়।
\item \textbf{রাতের আমল:} সপ্তাহে কোনো রাতে তাহাজ্জুদ + ১ রাকাত বিতর-অনুশীলন — সাথে পরিবারের কাউকে включ.
\item \textbf{জামাতের শিষ্টাচার:} জামাতে দাঁড়ানোর সঠিক স্থান (ডানে) ও অনুশীলনের নিয়ম — তাড়াহুড়া নয়, স্থিরতা।
\end{enumerate}
\medskip\hrule\medskip

\subsection*{৯. রেফারেন্স}

\begin{itemize}
\item সহীহ বুখারি — ১৩৮ (পূর্ণ রাতের ঘটনা)
\item সহীহ মুসলিম — ৭৬৩ (একই ঘটনার আভিধানিক বর্ণনা)
\item সুনানে আবু দাউদ — ১৩২৬-এর মতো তাহাজ্জুদ-বিতর সম্পর্কিত
\item কুরআন: আল-ইসরা ১৭:৭৯; আল-যারিয়াত ৫১:১৭-১৮
\end{itemize}

\end{document}
